% GENERATED FILE. Edit canonical lessons, then run npm run generate:books.
% canonical-source-hash: fnv1a64:7b7d5c281cc4a75d
% canonical-lessons: PT-C484-enxaqueca, PT-C484-tonto, PT-C484-tontura, PT-C484-nausea, PT-C484-vomitar

\chapter{Dores de cabeça e náuseas}
\label{ch:pt-dores-de-cabeca-e-nauseas}
\hlchaptermodality{484}

\begin{chapteropening}
\textbf{By the end of this chapter:} \emph{I can describe headaches, dizziness and nausea.}

Complete the last lesson of chapter 484: I can describe headaches, dizziness and nausea.
\end{chapteropening}

\section[a enxaqueca]{a enxaqueca — migraine}

\label{lesson:PT-C484-enxaqueca}



\begin{quote}
Before the new one: say the Portuguese for itch, then the Portuguese for to scratch an itch.
\end{quote}

\subsection*{You'll want to know: a enxaqueca}
\textbf{a enxaqueca} — \textquotedblleft{}migraine\textquotedblright{}.

\begin{quote}
Quando tenho enxaqueca, fico às escuras no quarto.
\end{quote}

\begin{quote}
\emph{When I have a migraine, I stay in my room in the dark.}
\end{quote}

Stronger than \emph{uma dor de cabeça}.

\begin{cousinweb}
\emph{enxaqueca} is usually traced to Arabic \emph{aš-šaqīqa}, a word for a headache on one side of the head.
\end{cousinweb}

\begin{grammarlens}[title={What you've built}]
One more word for telling, explaining and describing at B1.
\end{grammarlens}

\subsection*{Your turn}
\begin{itemize}
\raggedright
  \item \emph{Say it:} \emph{a enxaqueca}
  \item \emph{Say it:} \emph{a enxaqueca}, once more
  \item \emph{Say it:} a sentence of your own with \emph{a enxaqueca}
  \item \emph{Recall:} say the Portuguese for itch, then the Portuguese for to scratch an itch, then say \emph{a enxaqueca} again
\end{itemize}

\begin{tcolorbox}[breakable,colback=teal!4,colframe=teal!35!black,title={Before you move on}]
What does a enxaqueca mean? (\textquotedblleft{}Migraine\textquotedblright{}.) Say it once more.
\end{tcolorbox}
\section[tonto]{tonto — dizzy; silly}

\label{lesson:PT-C484-tonto}



\begin{quote}
Before the new one: say the Portuguese for to scratch an itch, then the Portuguese for migraine.
\end{quote}

\subsection*{You'll want to know: tonto}
\textbf{tonto} — \textquotedblleft{}dizzy; silly\textquotedblright{}.

\begin{quote}
Depois da viagem de barco, fiquei tonto.
\end{quote}

\begin{quote}
\emph{After the boat trip, I felt dizzy.}
\end{quote}

It also means silly: \emph{não sejas tonto}.

\begin{cousinweb}
\emph{tonto} is inherited from Latin \emph{attonitus} \textquotedblleft{}thunderstruck, stunned\textquotedblright{} (EtymDB).
\end{cousinweb}

\begin{grammarlens}[title={What you've built}]
One more word for telling, explaining and describing at B1.
\end{grammarlens}

\subsection*{Your turn}
\begin{itemize}
\raggedright
  \item \emph{Say it:} \emph{tonto}
  \item \emph{Say it:} \emph{tonto}, once more
  \item \emph{Say it:} a sentence of your own with \emph{tonto}
  \item \emph{Recall:} say the Portuguese for to scratch an itch, then the Portuguese for migraine, then say \emph{tonto} again
\end{itemize}

\begin{tcolorbox}[breakable,colback=teal!4,colframe=teal!35!black,title={Before you move on}]
What does tonto mean? (\textquotedblleft{}Dizzy; silly\textquotedblright{}.) Say it once more.
\end{tcolorbox}
\section[a tontura]{a tontura — dizziness}

\label{lesson:PT-C484-tontura}



\begin{quote}
Before the new one: say the Portuguese for migraine, then the Portuguese for dizzy.
\end{quote}

\subsection*{You'll want to know: a tontura}
\textbf{a tontura} — \textquotedblleft{}dizziness\textquotedblright{}.

\begin{quote}
Levantei-me depressa e senti uma tontura.
\end{quote}

\begin{quote}
\emph{I stood up quickly and felt dizzy.}
\end{quote}

Built on \emph{tonto}.

\begin{cousinweb}
\emph{tontura} is built on \emph{tonto} \textquotedblleft{}dizzy\textquotedblright{} with \emph{-ura} (EtymDB); \emph{tonto} is inherited from Latin \emph{attonitus} \textquotedblleft{}stunned\textquotedblright{}.
\end{cousinweb}

\begin{grammarlens}[title={What you've built}]
One more word for telling, explaining and describing at B1.
\end{grammarlens}

\subsection*{Your turn}
\begin{itemize}
\raggedright
  \item \emph{Say it:} \emph{a tontura}
  \item \emph{Say it:} \emph{a tontura}, once more
  \item \emph{Say it:} a sentence of your own with \emph{a tontura}
  \item \emph{Recall:} say the Portuguese for migraine, then the Portuguese for dizzy, then say \emph{a tontura} again
\end{itemize}

\begin{tcolorbox}[breakable,colback=teal!4,colframe=teal!35!black,title={Before you move on}]
What does a tontura mean? (\textquotedblleft{}Dizziness\textquotedblright{}.) Say it once more.
\end{tcolorbox}
\section[a náusea]{a náusea — nausea}

\label{lesson:PT-C484-nausea}



\begin{quote}
Before the new one: say the Portuguese for dizzy, then the Portuguese for dizziness.
\end{quote}

\subsection*{You'll want to know: a náusea}
\textbf{a náusea} — \textquotedblleft{}nausea\textquotedblright{}.

\begin{quote}
O cheiro do peixe deu-lhe náuseas.
\end{quote}

\begin{quote}
\emph{The smell of the fish made her feel sick.}
\end{quote}

Often plural: \emph{ter náuseas}.

\begin{cousinweb}
\emph{náusea} is a learned borrowing from Latin \emph{nausea} (EtymDB), from Greek \emph{nausía} \textquotedblleft{}seasickness\textquotedblright{}, built on \emph{naûs} \textquotedblleft{}ship\textquotedblright{}.
\end{cousinweb}

\begin{grammarlens}[title={What you've built}]
One more word for telling, explaining and describing at B1.
\end{grammarlens}

\subsection*{Your turn}
\begin{itemize}
\raggedright
  \item \emph{Say it:} \emph{a náusea}
  \item \emph{Say it:} \emph{a náusea}, once more
  \item \emph{Say it:} a sentence of your own with \emph{a náusea}
  \item \emph{Recall:} say the Portuguese for dizzy, then the Portuguese for dizziness, then say \emph{a náusea} again
\end{itemize}

\begin{tcolorbox}[breakable,colback=teal!4,colframe=teal!35!black,title={Before you move on}]
What does a náusea mean? (\textquotedblleft{}Nausea\textquotedblright{}.) Say it once more.
\end{tcolorbox}
\section[vomitar]{vomitar — to vomit, to be sick}

\label{lesson:PT-C484-vomitar}



\begin{quote}
Before the new one: say the Portuguese for dizziness, then the Portuguese for nausea.
\end{quote}

\subsection*{You'll want to know: vomitar}
\textbf{vomitar} — \textquotedblleft{}to vomit, to be sick\textquotedblright{}.

\begin{quote}
O miúdo vomitou no carro.
\end{quote}

\begin{quote}
\emph{The kid was sick in the car.}
\end{quote}

Everyday and neutral.

\begin{cousinweb}
\emph{vomitar} is a learned borrowing from Latin \emph{vomitāre} \textquotedblleft{}to vomit\textquotedblright{} (EtymDB), from \emph{vomere}.
\end{cousinweb}

\begin{grammarlens}[title={What you've built}]
One more word for telling, explaining and describing at B1.
\end{grammarlens}

\subsection*{Your turn}
\begin{itemize}
\raggedright
  \item \emph{Say it:} \emph{vomitar}
  \item \emph{Say it:} \emph{vomitar}, once more
  \item \emph{Say it:} a sentence of your own with \emph{vomitar}
  \item \emph{Recall:} say the Portuguese for dizziness, then the Portuguese for nausea, then say \emph{vomitar} again
\end{itemize}

\begin{tcolorbox}[breakable,colback=teal!4,colframe=teal!35!black,title={Before you move on}]
What does vomitar mean? (\textquotedblleft{}To vomit, to be sick\textquotedblright{}.) Say it once more.
\end{tcolorbox}
